% !TeX program = pdflatex
% !TeX encoding = UTF-8
% Object Counting Algebra -- the attached monolithic formulation.
% Source: Pasted text.txt. All fourteen sections retain their mathematical
% declarations and original order. Only the display layout is changed.
%
% Build (no shell escape, bibliography processor, or external files required):
%   pdflatex -interaction=nonstopmode -halt-on-error oca-defined.tex
%   pdflatex -interaction=nonstopmode -halt-on-error oca-defined.tex
%
% Page size remains A5. To change it, edit both occurrences of a5paper below.
% The original all-in-one gathered display cannot break across pages. Its
% outer alignment is therefore presented as unnumbered, breakable align*
% blocks. Nested cases, arrays, delimiters, and mathematical content remain.
\documentclass[10pt,a5paper]{article}

\usepackage[T1]{fontenc}
\usepackage[utf8]{inputenc}
\usepackage{lmodern}
\usepackage[a5paper,left=14mm,right=14mm,top=15mm,bottom=16mm]{geometry}
\usepackage{mathtools,amssymb,mathrsfs}
\usepackage{needspace}
\usepackage[hidelinks]{hyperref}
\hypersetup{
  pdftitle={Object Counting Algebra},
  pdfsubject={Monolithic mathematical formulation}
}

\allowdisplaybreaks[1]
\setlength{\jot}{4pt}
\setlength{\parindent}{0pt}
\setlength{\parskip}{6pt}
\setlength{\emergencystretch}{1em}
\raggedbottom

% Compact, unnumbered headings; no extra mathematical numbering is added.
\makeatletter
\renewcommand\section{\@startsection{section}{1}{\z@}%
  {2.0ex plus .6ex minus .2ex}{.8ex plus .2ex}%
  {\normalfont\large\bfseries}}
\makeatother

\begin{document}

{\LARGE\bfseries Object Counting Algebra\par}
\vspace{2pt}

The following is a single parameterized specification of the multiset construction in \emph{oca-discussion}. The general algebra, canonical encoding, operational semantics, and conditional specializations are incorporated into one formulation.

\[
\boxed{
\begin{aligned}
&\mathfrak O(A,T,\mathscr S,\sigma)
:=\Big\langle
\\[3pt]
&\quad
\mathcal C,\mathbf 0,\nu,\oplus,\star,\ominus,Q_{\mathrm R},Q_{\mathrm L};
\\[3pt]
&\quad
I,\rho,\delta,
\widehat{\oplus},\widehat{\star},\widehat{\ominus},
\widehat Q_{\mathrm R},\widehat Q_{\mathrm L};
\\[3pt]
&\quad
\mathscr S,\sigma,\sigma_I,\mathsf{Out}_I,\operatorname{Eval}
\Big\rangle .
\end{aligned}
}
\]

% Section 01 of 14
\Needspace{7\baselineskip}
\section*{Primitive domain}
\begin{align*}
&\mathbb N_0:=\{0,1,2,\ldots\},
\qquad
d\in\mathbb N_{\geq 1},
\qquad
[d]:=\{1,\ldots,d\},
\\
&A:=\{o_1,\ldots,o_d\},
\qquad
o_1<_A o_2<_A\cdots<_A o_d,
\\
&\mathcal C
:=
\left\{
\sum_{i=1}^{d}x_i o_i:
x_i\in\mathbb N_0
\right\}
\cong
\mathbb N_0^{d},
\\
&x=(x_1,\ldots,x_d)\in\mathcal C,
\qquad
(e_i)_k:=\mathbf 1_{\{i=k\}},
\\
&x=\bigoplus_{i=1}^{d}x_i e_i,
\qquad
\mathbf 0:=(0,\ldots,0),
\\
&x=y
\iff
\forall i\in[d],\ x_i=y_i,
\\
&x\preccurlyeq y
\iff
\forall i\in[d],\ x_i\leq y_i,
\\
&\nu:\mathcal C\longrightarrow\mathbb N_0,
\qquad
\nu(x):=\sum_{i=1}^{d}x_i,
\\
&\mathcal C_n:=\{x\in\mathcal C:\nu(x)=n\},
\qquad
|\mathcal C_n|=\binom{n+d-1}{d-1},
\\
&\mathcal C=\bigsqcup_{n\in\mathbb N_0}\mathcal C_n,
\qquad
|\mathcal C|=\aleph_0.
\end{align*}

% Section 02 of 14
\Needspace{7\baselineskip}
\section*{Canonical sequential indexing}
\begin{align*}
&x<_{\mathrm{lex}}y
\iff{}
\exists r\in[d]:
\left(
x_r<y_r
\ \land\
\forall h<r,\ x_h=y_h
\right),
\\
&x\prec y
\iff{}
\nu(x)<\nu(y)
\\*[-2pt]
&{}\lor
\left(
\nu(x)=\nu(y)
\ \land\
x<_{\mathrm{lex}}y
\right),
\\
&I:={}\mathbb N_0,
\\
&\rho:\mathcal C\longrightarrow I,
\qquad
\rho(x):={}
\left|\{y\in\mathcal C:y\prec x\}\right|,
\\
&p_i(x):={}\sum_{h=1}^{i-1}x_h,
\\
&\rho(x)={}
\binom{\nu(x)+d-1}{d}
\\*[-2pt]
&{}+
\sum_{i=1}^{d-1}
\sum_{a=0}^{x_i-1}
\binom{
\nu(x)-p_i(x)-a+d-i-1
}{
d-i-1
},
\\
&\sum_{\varnothing}:=0,
\qquad
\binom{n}{k}:=0
\quad(0\leq n<k),
\\
&\forall a\in I,\ \exists!x\in\mathcal C:\rho(x)=a,
\\
&\delta:=\rho^{-1}:I\longrightarrow\mathcal C,
\\
&\delta\circ\rho=\operatorname{id}_{\mathcal C},
\qquad
\rho\circ\delta=\operatorname{id}_{I},
\\
&\rho(\mathbf 0)=0,
\qquad
\rho(e_i)=d-i+1,
\\
&\rho(x)=\rho(y)\iff x=y.
\end{align*}

% Section 03 of 14
\Needspace{7\baselineskip}
\section*{Primitive composition and configuration operations}
\begin{align*}
&T=(T_{ij}^{k})_{i,j,k\in[d]}
\in\mathbb N_0^{d\times d\times d},
\\
&\mu_T:A\times A\longrightarrow\mathcal C,
\qquad
\mu_T(o_i,o_j):=\sum_{k=1}^{d}T_{ij}^{k}e_k,
\\
&\oplus,\star:\mathcal C\times\mathcal C
\longrightarrow\mathcal C,
\\
&(x\oplus y)_k:=x_k+y_k,
\\
&(x\star y)_k
:=
\sum_{i=1}^{d}\sum_{j=1}^{d}
x_i y_j T_{ij}^{k},
\\
&e_i\star e_j=\mu_T(o_i,o_j),
\\
&\ominus:\mathcal C\times\mathcal C
\rightharpoonup\mathcal C,
\\
&\operatorname{dom}(\ominus)
:=
\{(x,y)\in\mathcal C^2:y\preccurlyeq x\},
\\
&(x\ominus y)_k:=x_k-y_k
\qquad
\bigl((x,y)\in\operatorname{dom}(\ominus)\bigr),
\\
&\mathcal R(x,y)
:=
\{z\in\mathcal C:z\oplus y=x\}
\\*
&\hspace{39pt}=
\begin{cases}
\{x\ominus y\},&y\preccurlyeq x,\\
\varnothing,&y\not\preccurlyeq x,
\end{cases}
\\
&Q_{\mathrm R},Q_{\mathrm L}:
\mathcal C\times\mathcal C
\longrightarrow\mathcal P(\mathcal C),
\\
&Q_{\mathrm R}(x,y)
:=
\{q\in\mathcal C:q\star y=x\},
\\
&Q_{\mathrm L}(x,y)
:=
\{q\in\mathcal C:y\star q=x\}.
\end{align*}

% Section 04 of 14
\Needspace{7\baselineskip}
\section*{Unconditional algebraic laws}
\begin{align*}
&\forall x,y,z\in\mathcal C:
\\*
&(x\oplus y)\oplus z=x\oplus(y\oplus z),
\\
&x\oplus y=y\oplus x,
\qquad
x\oplus\mathbf 0=x,
\\
&x\oplus z=y\oplus z\iff x=y,
\\
&x\oplus y=\mathbf 0
\iff
x=y=\mathbf 0,
\\
&x\star(y\oplus z)
=(x\star y)\oplus(x\star z),
\\
&(x\oplus y)\star z
=(x\star z)\oplus(y\star z),
\\
&x\star\mathbf 0=\mathbf 0\star x=\mathbf 0,
\\
&\nu(x\oplus y)=\nu(x)+\nu(y),
\\
&y\preccurlyeq x
\Longrightarrow
(x\ominus y)\oplus y=x,
\\
&(x\oplus y)\ominus y=x,
\qquad
x\ominus x=\mathbf 0,
\\
&(\mathcal C,\oplus,\mathbf 0)
\cong
\operatorname{FreeCMon}(A),
\\
&\nu(x\star y)
=
\sum_{i=1}^{d}\sum_{j=1}^{d}
x_i y_j
\left(\sum_{k=1}^{d}T_{ij}^{k}\right),
\\
&\left[
\forall x,y\in\mathcal C,\
\nu(x\star y)=\nu(x)\nu(y)
\right]
\\*
&\hspace{35pt}\iff
\left[
\forall i,j\in[d],\
\sum_{k=1}^{d}T_{ij}^{k}=1
\right].
\end{align*}

% Section 05 of 14
\Needspace{7\baselineskip}
\section*{Associativity, commutativity, and multiplicative identity}
\begin{align*}
&\mathsf A_T
\iff{}
\forall i,j,k,m\in[d]:
\\*[-2pt]
&\sum_{\ell=1}^{d}
T_{ij}^{\ell}T_{\ell k}^{m}
=
\sum_{\ell=1}^{d}
T_{jk}^{\ell}T_{i\ell}^{m},
\\
&\mathsf C_T
\iff{}
\forall i,j,k\in[d],\
T_{ij}^{k}=T_{ji}^{k},
\\
&\mathcal U_T
:={}
\left\{
u\in\mathcal C:
\begin{array}{l}
\displaystyle
\forall i,k\in[d],\\[2pt]
\displaystyle
\sum_{j=1}^{d}u_jT_{ji}^{k}
=
\mathbf 1_{\{i=k\}}
=
\sum_{j=1}^{d}u_jT_{ij}^{k}
\end{array}
\right\},
\\
&\mathsf A_T
\iff{}
\forall x,y,z\in\mathcal C,\
(x\star y)\star z=x\star(y\star z),
\\
&\mathsf C_T
\iff{}
\forall x,y\in\mathcal C,\
x\star y=y\star x,
\\
&|\mathcal U_T|\leq1,
\\
&u\in\mathcal U_T
\Longrightarrow{}
u\neq\mathbf 0
\ \land\
\forall x\in\mathcal C,\
u\star x=x\star u=x,
\\
&\mathsf A_T\land u\in\mathcal U_T
\Longrightarrow{}
(\mathcal C,\oplus,\star,\mathbf 0,u)
\in\mathsf{Semiring}_1,
\\
&\mathsf A_T\land\mathsf C_T
\land u\in\mathcal U_T
\Longrightarrow{}
(\mathcal C,\oplus,\star,\mathbf 0,u)
\in\mathsf{CSemiring}_1,
\\
&u\in\mathcal U_T
\Longrightarrow{}
\forall n\in\mathbb N_0,\ \forall x\in\mathcal C:
\\*[-2pt]
&(nu)\star x=x\star(nu)=nx.
\end{align*}

% Section 06 of 14
\Needspace{7\baselineskip}
\section*{Division as nonnegative integer fibers}
\begin{align*}
&(M_{\mathrm R}(y))_{ki}
:=
\sum_{j=1}^{d}T_{ij}^{k}y_j,
\\
&(M_{\mathrm L}(y))_{ki}
:=
\sum_{j=1}^{d}T_{ji}^{k}y_j,
\\
&q\star y=M_{\mathrm R}(y)q,
\qquad
y\star q=M_{\mathrm L}(y)q,
\\
&\varepsilon\in\{\mathrm R,\mathrm L\},
\\
&Q_{\varepsilon}(x,y)
=
\{q\in\mathbb N_0^d:M_{\varepsilon}(y)q=x\},
\\
&Q_{\varepsilon}(x,\mathbf 0)
=
\begin{cases}
\mathcal C,&x=\mathbf 0,\\
\varnothing,&x\neq\mathbf 0,
\end{cases}
\\
&Q_{\varepsilon}(\mathbf 0,y)
=
\{q\in\mathbb N_0^d:M_{\varepsilon}(y)q=\mathbf 0\},
\\
&c_i^{\varepsilon}(y)
:=
\sum_{k=1}^{d}(M_{\varepsilon}(y))_{ki},
\\
&q\in Q_{\varepsilon}(x,y)
\Longrightarrow
\sum_{i=1}^{d}c_i^{\varepsilon}(y)q_i=\nu(x),
\\
&|Q_{\varepsilon}(x,y)|=\aleph_0
\\*[-2pt]
&\hspace{20pt}\iff
Q_{\varepsilon}(x,y)\neq\varnothing
\ \land\
\exists i\in[d],\ c_i^{\varepsilon}(y)=0,
\\
&\left[
\forall x\in\mathcal C,\
|Q_{\varepsilon}(x,y)|\leq1
\right]
\iff
\det M_{\varepsilon}(y)\neq0.
\end{align*}

% Section 07 of 14
\Needspace{7\baselineskip}
\section*{Transport to identifier arithmetic}
\begin{align*}
&\widehat{\oplus},\widehat{\star}:I^2\longrightarrow I,
\qquad
\widehat{\ominus}:I^2\rightharpoonup I,
\\
&\widehat Q_{\mathrm R},\widehat Q_{\mathrm L}:
I^2\longrightarrow\mathcal P(I),
\\
&a\widehat{\oplus}b
:=
\rho\bigl(\delta(a)\oplus\delta(b)\bigr),
\\
&a\widehat{\star}b
:=
\rho\bigl(\delta(a)\star\delta(b)\bigr),
\\
&\operatorname{dom}(\widehat{\ominus})
=
\{(a,b)\in I^2:\delta(b)\preccurlyeq\delta(a)\},
\\
&a\widehat{\ominus}b
:=
\rho\bigl(\delta(a)\ominus\delta(b)\bigr)
\quad
\bigl((a,b)\in\operatorname{dom}(\widehat{\ominus})\bigr),
\\
&\rho[U]:=\{\rho(x):x\in U\}
\qquad
(U\subseteq\mathcal C),
\\
&\widehat Q_{\varepsilon}(a,b)
:=
\rho\!\left[Q_{\varepsilon}(\delta(a),\delta(b))\right],
\\
&0_I:=\rho(\mathbf 0)=0,
\qquad
u\in\mathcal U_T\Longrightarrow 1_I:=\rho(u),
\\
&\nu_I:=\nu\circ\delta,
\\
&\nu_I(a\widehat{\oplus}b)
=
\nu_I(a)+\nu_I(b).
\end{align*}

% Section 08 of 14
\Needspace{7\baselineskip}
\section*{Typed outcomes and total evaluation}
\begin{align*}
&\mathsf{Out}_I
:=
\mathsf{Val}(I)
\sqcup
\mathsf{Sol}\bigl(\mathcal P(I)\bigr)
\sqcup
\{\mathsf{Fail}\},
\\
&\left|
\left\{
\mathsf{Val}(0),
\mathsf{Sol}(\varnothing),
\mathsf{Sol}(\{0\}),
\mathsf{Fail}
\right\}
\right|=4,
\\
&\Omega_{\mathrm{alg}}
:=
\{\oplus,\star,\ominus\},
\\
&\Omega
:=
\Omega_{\mathrm{alg}}
\sqcup
\{\operatorname{div}_{\mathrm R},
\operatorname{div}_{\mathrm L}\},
\\
&\operatorname{Eval}:
\Omega\times I\times I\longrightarrow\mathsf{Out}_I,
\\
&\operatorname{Eval}(\omega,a,b)
:=
\begin{cases}
\mathsf{Val}(a\widehat{\omega}b),
&
\begin{array}{l}
\omega\in\Omega_{\mathrm{alg}},\\
(a,b)\in\operatorname{dom}(\widehat\omega),
\end{array}
\\[5pt]
\mathsf{Fail},
&
\begin{array}{l}
\omega\in\Omega_{\mathrm{alg}},\\
(a,b)\notin\operatorname{dom}(\widehat\omega),
\end{array}
\\[5pt]
\mathsf{Sol}\bigl(\widehat Q_{\varepsilon}(a,b)\bigr),
&
\omega=\operatorname{div}_{\varepsilon}.
\end{cases}
\end{align*}

% Section 09 of 14
\Needspace{7\baselineskip}
\section*{Semantic interpretation and compatibility}
\begin{align*}
&\mathscr S:=(S,\oplus_S,\star_S,0_S),
\\
&\oplus_S,\star_S:S^2\longrightarrow S,
\qquad
0_S\in S,
\\
&\sigma:\mathcal C\longrightarrow S,
\qquad
\sigma_I:=\sigma\circ\delta,
\\
&\mathsf H_{\sigma}
\iff
\left\{
\begin{aligned}
&\sigma(\mathbf 0)=0_S,
\\
&\forall x,y\in\mathcal C:
\\
&\sigma(x\oplus y)
=
\sigma(x)\oplus_S\sigma(y),
\\
&\sigma(x\star y)
=
\sigma(x)\star_S\sigma(y),
\end{aligned}
\right.
\\
&\mathcal R_S(s,t)
:=
\{v\in S:v\oplus_S t=s\},
\\
&Q_{\mathrm R}^{S}(s,t)
:=
\{v\in S:v\star_S t=s\},
\\
&Q_{\mathrm L}^{S}(s,t)
:=
\{v\in S:t\star_S v=s\},
\\
&\sigma[U]:=\{\sigma(x):x\in U\},
\\
&\mathsf H_{\sigma}
\Longrightarrow
\sigma[\mathcal R(x,y)]
\subseteq
\mathcal R_S(\sigma(x),\sigma(y)),
\\
&\mathsf H_{\sigma}
\Longrightarrow
\sigma[Q_{\mathrm R}(x,y)]
\subseteq
Q_{\mathrm R}^{S}(\sigma(x),\sigma(y)),
\\
&\mathsf H_{\sigma}
\Longrightarrow
\sigma[Q_{\mathrm L}(x,y)]
\subseteq
Q_{\mathrm L}^{S}(\sigma(x),\sigma(y)),
\\
&x\sim_{\sigma}y
\iff
\sigma(x)=\sigma(y),
\\
&\mathsf H_{\sigma}
\Longrightarrow
\left[
\begin{array}{c}
x\sim_{\sigma}x',
\quad y\sim_{\sigma}y'
\\[2pt]
\Downarrow
\\[2pt]
x\oplus y\sim_{\sigma}x'\oplus y',
\\
x\star y\sim_{\sigma}x'\star y'
\end{array}
\right],
\\
&\left[
\begin{array}{l}
\forall x,x',y,y'\in\mathcal C:
\\
x\sim_{\sigma}x'\ \land\ y\sim_{\sigma}y'
\\
\qquad\Longrightarrow
\bigl(y\preccurlyeq x
\iff y'\preccurlyeq x'\bigr)
\end{array}
\right]
\\*[-2pt]
&\hspace{40pt}\Longrightarrow
\left[
x\sim_{\sigma}y\iff x=y
\right].
\end{align*}

% Section 10 of 14
\Needspace{7\baselineskip}
\section*{Reindexing invariance}
\begin{align*}
&\pi:I\overset{\sim}{\longrightarrow}I,
\qquad
\rho_{\pi}:=\pi\circ\rho,
\qquad
\delta_{\pi}:=\delta\circ\pi^{-1},
\\
&\widehat\omega_{\pi}
:=
\pi\circ\widehat\omega
\circ(\pi^{-1}\times\pi^{-1})
\qquad
(\omega\in\Omega_{\mathrm{alg}}),
\\
&\widehat Q_{\varepsilon,\pi}(a,b)
:=
\pi\!\left[
\widehat Q_{\varepsilon}(\pi^{-1}(a),\pi^{-1}(b))
\right],
\\
&\sigma_{I,\pi}:=\sigma_I\circ\pi^{-1},
\\
&0_{I,\pi}=\pi(0_I),
\qquad
u\in\mathcal U_T\Longrightarrow
1_{I,\pi}=\pi(1_I),
\\
&\mathfrak O_I\cong_{\pi}\mathfrak O_{I,\pi}.
\end{align*}

% Section 11 of 14
\Needspace{7\baselineskip}
\section*{Finite monoid specialization}
\begin{align*}
&m:[d]\times[d]\longrightarrow[d],
\qquad
h\in[d],
\\
&\mathsf M_{m,h}(T)
\iff
\left\{
\begin{aligned}
&T_{ij}^{k}
=
\mathbf1_{\{m(i,j)=k\}}
&&\forall i,j,k\in[d],
\\
&m(m(i,j),k)=m(i,m(j,k))
&&\forall i,j,k\in[d],
\\
&m(h,i)=m(i,h)=i
&&\forall i\in[d],
\end{aligned}
\right.
\\
&\mathsf M_{m,h}(T)
\Longrightarrow
M:=([d],m,h)\in\mathsf{Mon},
\\
&\mathsf M_{m,h}(T)
\Longrightarrow
(\mathcal C,\oplus,\star,\mathbf0,e_h)
\cong
\mathbb N_0[M],
\\
&\mathsf M_{m,h}(T)
\Longrightarrow
\mathcal U_T=\{e_h\},
\qquad
1_I=\rho(e_h)=d-h+1,
\\
&\mathsf M_{m,h}(T)
\Longrightarrow
\nu(x\star y)=\nu(x)\nu(y),
\\
&\mathsf M_{m,h}(T)
\Longrightarrow
\nu:\mathbb N_0[M]\longrightarrow\mathbb N_0
\text{ is a unital semiring homomorphism},
\\
&\mathsf M_{m,h}(T)
\Longrightarrow
\nu_I(a\widehat\star b)=\nu_I(a)\nu_I(b),
\\
&\mathsf M_{m,h}(T)\land y\neq\mathbf0
\\*[-2pt]
&\hspace{20pt}\Longrightarrow
Q_{\varepsilon}(x,y)
\subseteq
\{q\in\mathcal C:\nu(q)\nu(y)=\nu(x)\},
\\
&\mathsf M_{m,h}(T)\land y\neq\mathbf0
\land \nu(y)\nmid\nu(x)
\\*[-2pt]
&\hspace{20pt}\Longrightarrow
Q_{\varepsilon}(x,y)=\varnothing,
\\
&\mathsf M_{m,h}(T)\land y\neq\mathbf0
\land N:=\frac{\nu(x)}{\nu(y)}\in\mathbb N_0
\\*[-2pt]
&\hspace{20pt}\Longrightarrow
|Q_{\varepsilon}(x,y)|
\leq
\binom{N+d-1}{d-1}<\infty,
\\
&\mathsf M_{m,h}(T)\land y\neq\mathbf0
\Longrightarrow
Q_{\varepsilon}(\mathbf0,y)=\{\mathbf0\}.
\end{align*}

% Section 12 of 14
\Needspace{7\baselineskip}
\section*{Admissibility restriction}
\begin{align*}
&\mathbf0\in D\subseteq\mathcal C,
\qquad
\omega\in\{\oplus,\star,\ominus\},
\\
&\operatorname{dom}(\omega_D)
:=
\left\{
(x,y)\in D^2:
\begin{array}{l}
(x,y)\in\operatorname{dom}(\omega),\\
\omega(x,y)\in D
\end{array}
\right\},
\\
&\omega_D(x,y):=\omega(x,y)
\qquad
\bigl((x,y)\in\operatorname{dom}(\omega_D)\bigr),
\\
&Q_{\varepsilon,D}(x,y)
:=
Q_{\varepsilon}(x,y)\cap D,
\qquad
(x,y)\in D^2,
\\
&\rho_D(x)
:=
|\{y\in D:y\prec x\}|,
\\
&I_D:=\rho_D[D]
=
\begin{cases}
\{0,\ldots,|D|-1\},&|D|<\infty,\\
\mathbb N_0,&|D|=\aleph_0,
\end{cases}
\\
&\delta_D:=\rho_D^{-1},
\qquad
\sigma_D:=\sigma|_D,
\\
&\operatorname{dom}(\oplus_D)=D^2
\iff
\oplus[D^2]\subseteq D,
\\
&\operatorname{dom}(\star_D)=D^2
\iff
\star[D^2]\subseteq D,
\\
&\widehat\omega_D
:=
\rho_D\circ\omega_D\circ(\delta_D\times\delta_D),
\\
&\widehat Q_{\varepsilon,D}(a,b)
:=
\rho_D\!\left[
Q_{\varepsilon,D}(\delta_D(a),\delta_D(b))
\right].
\end{align*}

% Section 13 of 14
\Needspace{7\baselineskip}
\section*{Signed extension}
\begin{align*}
&\mathcal C^{\pm}:=\mathbb Z^d
\cong
\operatorname{Gr}(\mathcal C,\oplus),
\\
&(x,y)\sim_{\mathrm{Gr}}(x',y')
\iff
x\oplus y'=x'\oplus y,
\\
&(\mathcal C\times\mathcal C)/{\sim_{\mathrm{Gr}}}
\longrightarrow
\mathbb Z^d,
\qquad
[(x,y)]\longmapsto x-y,
\\
&(z\oplus_{\pm}w)_k:=z_k+w_k,
\qquad
(z\ominus_{\pm}w)_k:=z_k-w_k,
\\
&(z\star_{\pm}w)_k
:=
\sum_{i=1}^{d}\sum_{j=1}^{d}
z_iw_jT_{ij}^{k},
\\
&\ominus_{\pm}:
(\mathcal C^{\pm})^2\longrightarrow\mathcal C^{\pm},
\\
&Q_{\mathrm R}^{\pm}(z,w)
:=
\{q\in\mathbb Z^d:q\star_{\pm}w=z\},
\\
&Q_{\mathrm L}^{\pm}(z,w)
:=
\{q\in\mathbb Z^d:w\star_{\pm}q=z\},
\\
&\|z\|_1:=\sum_{i=1}^{d}|z_i|,
\\
&z\prec_{\pm}w
\iff
\|z\|_1<\|w\|_1
\\*[-2pt]
&\hspace{71pt}\lor
\left(
\|z\|_1=\|w\|_1
\ \land\
z<_{\mathrm{lex}}w
\right),
\\
&\rho_{\pm}(z)
:=
|\{w\in\mathbb Z^d:w\prec_{\pm}z\}|,
\\
&\rho_{\pm}:\mathbb Z^d
\overset{\sim}{\longrightarrow}\mathbb N_0,
\qquad
\delta_{\pm}:=\rho_{\pm}^{-1},
\\
&\mathsf A_T\land u\in\mathcal U_T
\\*[-2pt]
&\hspace{20pt}\Longrightarrow
(\mathbb Z^d,\oplus_{\pm},\star_{\pm},\mathbf0,u)
\in\mathsf{Ring}_1,
\\
&\mathsf M_{m,h}(T)
\\*[-2pt]
&\hspace{20pt}\Longrightarrow
(\mathbb Z^d,\oplus_{\pm},\star_{\pm},\mathbf0,e_h)
\cong
\mathbb Z[M].
\end{align*}

% Section 14 of 14
\Needspace{7\baselineskip}
\section*{Neutral–toggle instance}
\begin{align*}
&d=2,
\qquad
o_1=\mathrm{neutral},
\qquad
o_2=\mathrm{toggle},
\\
&T_{11}^{1}=T_{12}^{2}=T_{21}^{2}=T_{22}^{1}=1,
\\
&T_{ij}^{k}=0
\quad
\text{for all remaining }(i,j,k),
\\
&(p,q)\star(r,s)
=
(pr+qs,\ ps+qr),
\\
&\rho(p,q)
=
\binom{p+q+1}{2}+p,
\\
&\begin{array}{c|cccccc}
a&0&1&2&3&4&5\\ \hline
\delta(a)&(0,0)&(0,1)&(1,0)&(0,2)&(1,1)&(2,0)
\end{array}
\\
&1\widehat{\oplus}2=4,
\qquad
4\widehat{\ominus}1=2,
\qquad
1\widehat{\star}1=2,
\\
&\widehat Q_{\mathrm R}(4,4)
=
\widehat Q_{\mathrm L}(4,4)
=
\{1,2\},
\\
&Q_{\mathrm R}\bigl((1,1),(1,1)\bigr)
=
\{(0,1),(1,0)\},
\\
&Q_{\mathrm R}^{\pm}\bigl((1,1),(1,1)\bigr)
=
\{(t,1-t):t\in\mathbb Z\}.
\end{align*}

\end{document}
